\printedpage{44}
% head: A. Grothendieck : Dérivateurs, Ch. XIII, Édition des Universités de Montpellier II et Paris VII

\begin{itemize}
\item[c)] \textit{$I'$ large relative to $\pi_0$.}
\end{itemize}

\texttt{[The lemma as stated is obvious for the same reasons as above (even without the hypothesis of the footnote). The proof that follows in fact proves a finer result:}

\texttt{\textsc{Lemma} 1$'$: \textit{Let $\pi_0$, $\pi$ be two cardinals, with $\pi_0<\pi$ and $\pi^{\pi_0}=\pi$, $I$ an ordered set large relative to $\pi_0$, and $I'_0$ a subset of $I$ of cardinality $\le\pi$. Then there exists a subset $I'$ of $I$ satisfying the conditions:}}

\begin{itemize}
\item[\texttt{a)}] $I'\supset I'_0$\texttt{.}
\item[\texttt{b)}] $\card I'\le\pi$\texttt{.}
\item[\texttt{c)}] \texttt{\textit{$I'$ large relative to $\pi_0$. ]}}
\end{itemize}

We are going to try to construct $I'$ by transfinite induction, defining the $I'_\alpha\subset I$ as follows:
\[
\begin{aligned}
I'_0 &= I'_0\;, \\
I'_\alpha &= \textstyle\bigcup_{\beta<\alpha} I'_\beta \quad \text{if $\alpha$ is a limit ordinal,} \\
I'_{\alpha+1} &= \left\{\parbox{0.6\textwidth}{derived from $I'_\alpha$ as follows: for every subset $J$ of $I'_\alpha$ such that $\card J\le\pi_0$, which is not bounded above in $I'_\alpha$, we choose an upper bound $i_J$ in $I$ (\texttt{[which]} exists, since $I$ \texttt{[is]} large relative to $\pi$, hence to $\pi_0$) and $I'_{\alpha+1}$ is obtained by adjoining these $i_J$ to $I'_\alpha$.}\right.
\end{aligned}
\]

\msp{page 81}
Let then $\alpha_0$ be the first ordinal of cardinality $>\pi_0$, and consider
\[
I' \;=\; I'_{\alpha_0} \;=\; \bigcup\limits_{\substack{\alpha<\alpha_0\\ \textit{i.e.}\ \card\alpha\le\pi_0}} I'_\alpha .
\]
Then $I'$ contains $I'_0$, it is large relative to $\pi_0$, but what is its cardinality? Let us prove by transfinite induction that
\[
\card I'_\alpha \;\le\pi \qquad\quad \text{if } \alpha<\alpha_0,\ \ \textit{i.e.}\ \ \card\alpha\le\pi_0 .
\]
This is O.K. for $\alpha=0$ by the hypothesis on $I'_0$. If it is true for $I'_\alpha$, is it true for $I'_{\alpha+1}$? We adjoin to $I'_\alpha$ a set of elements $i_J$ whose cardinality is $\le\card\mathfrak{P}_{\le\pi_0}(I'_\alpha)$, hence of cardinality $\le\card(I'_\alpha)^{\pi_0}\le\pi^{\pi_0}$, and if we suppose
\[
\pi^{\pi_0} \;=\; \pi,
\]
we win. Finally suppose that $\alpha$ is a limit ordinal (of cardinality $\le\pi_0$). So we have
\[
I'_\alpha \;=\; \bigcup_{\beta<\alpha} I'_\beta ,
\]
so
\printedpage{45}
% head: A. Grothendieck : Dérivateurs, Ch. XIII, Édition des Universités de Montpellier II et Paris VII
\msp{page 82}
\[
\card I'_\alpha \;\le\; \underbrace{(\card\alpha)\pi}_{\le\pi\pi=\pi^2=\pi} \qquad\quad \text{since the } \card I'_\beta\le\pi,
\]
we win again.

Finally we find, by the same argument,
\[
\card I'_{\alpha_0} \;\le\; \underbrace{(\card\alpha_0)\pi\le\pi^2}_{\substack{\text{if }\card\alpha_0\le\pi,\\ \textit{i.e.}\text{ if }\pi_0<\pi}} = \pi,
\]
so we win again.

Thus, we have found a second direct system
\[
(X_{\tilde\imath})_{\tilde\imath\in\tilde I}
\]
of the same type as the initial system, having as direct limit $X$ (that it has the same direct limit is immediate). But what have we gained? This: in $\tilde I$ every subset of cardinality $\le\pi$ and large relative to $\pi_0$ is not only \textit{bounded above} in $\tilde I$, but even admits a least upper bound there. This follows from % margin: double vertical rule beside the four lines "inductive est immédiat" ... "supérieure. Cela résulte du"

\msp{page 83}
\textbf{Lemma 2:} \textit{Let $I$ be large relative to $\pi>\pi_0$, with $\pi^{\pi_0}=\pi$} \texttt{[it suffices that $\pi\ge\pi_0$]} \textit{(cf. lemma 1). Let $(I'_\alpha)_{\alpha\in A}$ be a set of subsets of $I$ which are in $\tilde I$,} i.e.

\begin{itemize}
\item[a)] \textit{of cardinality $\le\pi$, and}
\item[b)] \textit{large relative to $\pi_0$,}
\end{itemize}

\textit{and suppose that} $\card A\le\pi$ \textit{and that $A$, ordered by inclusion, is large relative to $\pi_0$. Then $I'=\bigcup_{\alpha\in A} I'_\alpha$ is still in $\tilde I$,} i.e. \textit{$I'$ is of cardinality $\le\pi$ (since of cardinality $\le\pi^2=\pi$), and moreover, large relative to $\pi_0$.} (Immediate verification.)

Thus, we have proved this:

% margin: single vertical rule beside Proposition 1, from here to the end of the page
\textbf{Proposition 1.} \textit{Let $E$ be} \texttt{[a]} \textit{big category, with cardinal filtration} $(\Filt^\pi)_{\pi\ge\pi_0}$. \textit{Let $\pi$ be a cardinal with}
\[
\pi>\pi_0, \qquad\quad \pi^{\pi_0}=\pi \qquad\quad \textit{(e.g. $\pi=2^c$, with $c\ge\pi_0$)} \quad \texttt{[}\pi\ge\pi_0\ \texttt{suffices].}
\]
\textit{Then for every $X$ in $E$, we can find a direct system $(X_i)_{i\in I}$, with limit $X$, with}

\begin{itemize}
\item[a)] \[ \left\{ \begin{aligned} & X_i\in\Filt^\pi \ \ \textit{for every $i\in I$,} \\ & \ I \ \textit{large relative to $\pi$,} \end{aligned} \right. \]
\end{itemize}

\textit{satisfying the following additional conditions:}

\begin{itemize}
\item[b)] \textit{Every subset $J$ of $I$ of cardinality $\le\pi$, and large relative to $\pi_0$, admits a least upper bound $i_J$ in $I$.}
\item[c)] \textit{We have, under these conditions}
\end{itemize}

\printedpage{46}
% head: A. Grothendieck : Dérivateurs, Ch. XIII, Édition des Universités de Montpellier II et Paris VII
\msp{page 84}
\[
X_{i_J} \xleftarrow{\ \sim\ } \varinjlim_{i\in J} X_i .
\]
\textit{Moreover, if $X\in\Filt^{\pi'}$ with $\pi'>\pi$, we can take $I$ of cardinality $\le\pi'^{\pi}$ (hence of cardinality $\le\pi'$ if $\pi'^{\pi}=\pi'$, e.g. if $\pi'=2^{c'}$, with $c'\ge\pi$).}

\textbf{Definition.} We call a direct system $(X_i)_{i\in I}$, satisfying conditions \textit{a}), \textit{b}), \textit{c}) above, a \textit{well-adapted direct system} (to the given cardinal filtration, for the given cardinals $\pi_0$, $\pi$), and the corresponding representation
\[
X \xleftarrow{\ \sim\ } \varinjlim X_i
\]
a \textit{well-adapted representation} of $X$ in terms of objects of $\Filt^\pi$.

\textbf{NB} When $E$ is stable under small filtering $\varinjlim$, the situation simplifies considerably. We can then take
\[
\tilde I \;=\; \{I'\subset I \mid \card I'\le\pi,\ I' \text{ filtering}\}
\]
(without burdening ourselves with the condition $I'$ large relative to $\pi_0$) and again define the direct system $(X_{I'})_{I'\in\tilde I}$ by
\[
X_{I'} \;=\; \varinjlim_{i\in I'} X_i \;.
\]
We again have
\msp{page 85}
\[
\card\tilde I \;\le\; (\card I)^\pi \;\le\; \pi'^{\pi} \qquad\quad \text{if }\ \card I\le\pi'^{\pi}
\]
(in particular, if $X\in\Filt^{\pi'}$, and we take $I$ as indicated in condition $c$ of cardinal filtrations). Moreover, this time, the fact that $\tilde I$ is large relative to $\pi$ is trivial, we even have better, apart from $I$ being filtering:

\begin{itemize}
\item[$(*)$] \textit{Every} subset $(I'_\alpha)_{\alpha\in A}$ of $\tilde I$ which is filtering and of cardinality $\le\pi$ admits a least upper bound $\tilde\imath_A$,
\end{itemize}

namely the union (which is indeed filtering). Moreover, the direct system $(X_{\tilde\imath})_{\tilde\imath\in\tilde I}$ has the stronger property that
\[
X_{\tilde\imath_A} \xleftarrow{\ \sim\ } \varprojlim_{\tilde\imath\in A} X_{\tilde\imath} \;. % [sic] \varprojlim (arrow under lim pointing left) -> \varinjlim, inductive limit of the inductive system X_ĩ over the filtrant part A
\]
Thus we have

\textbf{Proposition 1 bis.} \textit{Suppose $E$ is stable under small filtering $\varinjlim$, and} $\Filt^\pi$ \textit{stable under these, where $\pi\ge\pi_0$ is given. Then for every $X$ in $E$, we can find a direct system $(X_i)_{i\in I}$ with direct limit $X$, with}

\begin{itemize}
\item[a)] $X_i\in\Filt^\pi$ \textit{for every $i\in I$, $I$ an ordered set large relative to $\pi$.}
\item[b)] \textit{Every filtering subset $J$ of $I$ of cardinality $\le\pi$ admits a least upper bound $i_J$ in $I$.}
\item[c)] \textit{Under these conditions}
\end{itemize}

\printedpage{47}
% head: A. Grothendieck : Dérivateurs, Ch. XIII, Édition des Universités de Montpellier II et Paris VII
\msp{page 86}
\[
X_{i_J} \;=\; \varinjlim_{i\in J} X_i \;.
\]
Such direct systems will be called \textit{``perfectly adapted''} (to the given cardinal filtration, for the given cardinal $\pi$ – $\pi_0$ no longer figures in it, except as the beginning of the cardinal filtration and nothing more).

Since in practice, we will have to handle mostly $E$ stable under small filtering $\varinjlim$, and cardinal filtrations stable accordingly (strengthening of condition \textit{b}) of cardinal filtrations), it is this case, and the notion of \textit{perfectly adapted representation} of $X$ as $\varinjlim\limits_{I} X_i$ ($X_i\in\Filt^\pi$), that I will put forward. The ``general'' case, where we suppose only \texttt{[that]} $E$ satisfies $L_{\pi_0}$, is more convoluted and is handled less prettily (\footnote[41]{We will call a cardinal filtration \textit{strict,} if the $\underline{\Filt}^\pi$ are stable under filtering $\varinjlim$ of cardinality $\le\pi$ (and these small filtering limits exist in $E$).}).

\textbf{Definition 2.} Let $E$ be equipped with a strict cardinal filtration $(\underline{\Filt}^\pi)_{\pi\ge\pi_0}$. Let $M$ be a subset of $\Ob E$, $\underline{M}$ the full subcategory of $E$ which it defines, and let $\pi$ be a cardinal $\ge\pi_0$. We say that $M$ is \textit{$\pi$-accessible,} or that $\underline{M}$ is \textit{$\pi$-accessible}, % [sic] où que -> ou que, disjunction "ou" not the relative "où"
\msp{page 87}
if the following conditions are satisfied. For every direct system
\[
(X_i)_{i\in I}
\]
perfectly adapted to the cardinal filtration for $\pi$ (hence $I$ a filtering ordered set such that every filtering subset $J$ of cardinality $\le\pi$ has a least upper bound $i_J$, and $X_{i_J}\xleftarrow{\ \sim\ }\varinjlim\limits_{j\in J} X_j$); let $I_M$ be the set of $i\in I$ such that $X_i\in M$. With this understood, we want this: for $X$ ($=\varinjlim X_i$) to be in $M$, it is necessary and sufficient that $I_M$ be cofinal in $I$ (so that $X$ appears as a filtering $\varinjlim$, with index set $I_M$ large relative to $\pi$, of elements of $M\cap\underline{\Filt}^\pi\overset{\text{def}}{=}M^\pi$).

\msp{page 88}
I feel like taking up again the definition of the $\pi$-accessibility of $M$, for a given cardinal filtration, by way of equivalent conditions. For every direct system $(X_i)_{i\in I}$ in $M$ \texttt{[rather in $E$]}, I keep the notation
\[
I_M \;=\; \{i\in I \mid X_i\in M\}.
\]
I set, for every cardinal $\pi\ge\pi_0$,
\[
M^\pi \;=\; M\cap\Filt^\pi .
\]

\begin{itemize}
\item[\circled{a$_\pi$}] $M$ is the set of objects of the form
\[
\varinjlim_{i\in I} X_i \;,
\]
where the $X_i$ are in $M^\pi$, and $I$ is large relative to $\pi$.
\end{itemize}

\printedpage{48}
% head: A. Grothendieck : Dérivateurs, Ch. XIII, Édition des Universités de Montpellier II et Paris VII

\begin{itemize}
\item[\circled{b$_\pi$}] Let $(X_i)_{i\in I}$ be a direct system in $M$ \texttt{[rather in $E$]} perfectly adapted to the cardinal filtration for $\pi$, \textit{i.e.} $I$ filtering and every filtering subset $J$ of $I$ of cardinality $\le\pi$ admits a least upper bound $i_J$, and $X_{i_J}\xleftarrow{\ \sim\ }\varinjlim\limits_{J} X_i$, and finally $X_i\in\Filt^\pi$ for every $i\in I$. Let $X=\varinjlim X_i$. In order to have $X\in M$, it is necessary and sufficient that $I_M$ be cofinal in $I$.
\item[\circled{c$_\pi$}] $\underline{M}^\pi=\underline{M}\cap\underline{\Filt}^\pi$ is a subcategory of $\underline{M}$ generating by strict epimorphisms, $\underline{M}$ is stable under $\varinjlim$ large relative to $\pi$, finally for every $X\in M$, $\underline{M}^\pi/X$ is a filtering category.
\end{itemize}

\msp{page 89}
I presume that these conditions are equivalent.

$b_\pi\Longrightarrow a_\pi$. We suppose \textit{b}) holds, we want to prove \textit{a}). First that every $X\in M$ is of the stated form $\varinjlim\limits_{I} X_i$, with the $X_i\in M^\pi$ and $I$ large relative to $\pi$. By \texttt{[the]} proposition 1, we know that we have a perfect adapted presentation for $\pi$, $X=\varinjlim\limits_{i\in I} X_i$. By \textit{b}) ``it is necessary'', the subset $I_M$ of $I$ is cofinal in $I$, so since $I$ is large relative to $\pi$, so is $I_M$, and we will have (setting $J=I_M$)
\[
X \;=\; \varinjlim_{i\in J} X_i, \qquad\quad X_i\in M^\pi,\ \ J \text{ large relative to } \pi,
\]
O.K. Conversely, let us prove that if $X$ can be put in the form stated in \textit{a}), it is in $M$. This will follow from

\textbf{Lemma 1:} \textit{If condition} b) (\textit{or} a)) \textit{is satisfied, $M$ is stable under direct limits large relative to $\pi$.}

We must prove that if
\msp{page 90}
\[
X \;=\; \varinjlim_{i\in I} X_i \;,
\]
with $I$ large relative to $\pi$, the $X_i$ in $M$, then $X\in M$. Condition \textit{b}) ``it is sufficient'' ensures this when we suppose in addition the $X_i$ in $\Filt^\pi$ (hence in $M^\pi$), so that we will have $I_M=I$, hence $I_M$ cofinal, hence $X\in M$ (\footnote[42]{Not quite – we assume more there, about the direct system. We cannot avoid, even in the case $X_i\in M^\pi$ $\forall\, i\in I$, making the argument of the following pages.}). If we do not suppose the $X_i$ in $\Filt^\pi$, we must work a little more.

Go straight to the part of the proof below, page 96, which proves \textit{b}) $\Longrightarrow$ \textit{a}). The rest of the proof will prove the

\textbf{Lemma 2.} \textit{If every element of Filt$^\pi$ is $\pi$-accessible, then} b$_\pi$) \textit{(and even} a$_\pi$\textit{) implies that $M$ is stable under $\varinjlim$ large relative to $\pi$.}